% ---------------------------------------------------------------------------
% KAUST Academy lab report — TRELLIS.2 print-prep project
%
% This file is self-contained and compiles on its own (pdflatex + biber, or
% just pdflatex if you switch to plain BibTeX). If you have the official
% KAUST Academy .cls/.sty from the template, you can instead paste the body
% (everything between \section{Introduction} and \printbibliography) into it
% and drop this preamble.
% ---------------------------------------------------------------------------
\documentclass[11pt,a4paper]{article}

\usepackage[T1]{fontenc}
\usepackage[utf8]{inputenc}
\usepackage[margin=2.5cm]{geometry}
\usepackage{graphicx}
\usepackage{booktabs}
\usepackage{amsmath}
\usepackage{xcolor}
\usepackage{listings}
\usepackage{enumitem}
\usepackage{tcolorbox}

% KAUST / KAUST Academy palette (approximated; replace with the official
% definitions if you are compiling inside the real template). Defined before
% hyperref so the link colour can reference kateal.
\definecolor{kanavy}{HTML}{15355E}
\definecolor{kateal}{HTML}{00A0AF}
\definecolor{kagold}{HTML}{E8A33D}
\definecolor{kaorange}{HTML}{E2725B}
\definecolor{kalime}{HTML}{8CBF3F}
\definecolor{kagray}{HTML}{6E7B85}

\usepackage[colorlinks=true,allcolors=kateal]{hyperref}
\usepackage[backend=bibtex,style=numeric,sorting=none]{biblatex}
\addbibresource{bibliography.bib}

\lstset{
  basicstyle=\ttfamily\footnotesize,
  keywordstyle=\color{kanavy}\bfseries,
  commentstyle=\color{kagray}\itshape,
  stringstyle=\color{kalime},
  breaklines=true,
  frame=none,
  backgroundcolor=\color{black!3},
  captionpos=b,
}

\title{\vspace{-1.5cm}\color{kanavy}\bfseries From One Photo to a Printable Object:\\
A Validated Image-to-3D Pipeline on Apple Silicon}

\author{%
  Student One, Student Two, Student Three,\\
  Student Four, Student Five, Student Six,\\
  Mentor Name\\[0.3em]
  \small\texttt{student.one@kaust.edu.sa, student.two@kaust.edu.sa,}\\
  \small\texttt{student.three@kaust.edu.sa, mentor@kaust.edu.sa}
}
\date{July 16, 2026}

\begin{document}
\maketitle

\begin{tcolorbox}[colframe=kanavy,colback=white,arc=3mm,title={\bfseries\color{white}Abstract}]
We present a single-image-to-3D-print pipeline that runs entirely on Apple
Silicon and, unlike existing image-to-3D tools, \emph{measures} what it does to
the geometry it repairs. Generative image-to-3D models such as TRELLIS.2
produce meshes that look correct in a viewer but are not manufacturable: our
test asset, generated from one photograph, contains \textbf{2.2M triangles
spread over 4{,}518 disconnected components} and is not watertight, so a slicer
cannot print it. Repairing such a mesh is inherently destructive --- the fix
that guarantees watertightness (voxel remeshing) also silently reshapes the
object --- and no existing tool reports how much shape was lost. Our system has
three parts: (i) a port of the 4B-parameter TRELLIS.2 model from CUDA to
PyTorch~MPS/Metal, generating a textured mesh from one image in about five
minutes on an M4~Pro; (ii) a print-preparation stage that discards
floating-point noise, guarantees watertightness, and flood-fills the interior
solid; and (iii) a \textbf{validation layer} --- the core contribution --- that
quantifies the damage using topology-agnostic symmetric Chamfer and Hausdorff
surface distances, volume change, and printability diagnostics. Because repair
changes vertex count by over an order of magnitude, correspondence-based
metrics are unusable; our point-to-surface formulation requires only that the
two meshes occupy the same space. The validation layer turns an unverifiable
black box into a measured one: it reports exactly \textbf{0.00\% deviation}
when post-processing is a no-op, and on the real asset it quantifies the
accuracy/manufacturability trade-off that the user is being asked to accept.
\end{tcolorbox}

\vspace{0.5em}
\noindent\textbf{Date:} July 16, 2026\\
\textbf{Code:} \url{https://github.com/shivampkumar/trellis-mac}

\tableofcontents
\newpage

% ===========================================================================
\section{Introduction}
% ===========================================================================

\paragraph{Motivation.}
Digital fabrication labs receive a recurring request: \emph{``I have a photo of
this object --- can you print me one?''} Until recently the answer required a
trained user and hours of manual CAD work. Generative image-to-3D models have
made the first half of that request nearly trivial: TRELLIS.2
\cite{xiang2025trellis2} reconstructs a textured 3D asset from a single
photograph in minutes. The promise is significant for teaching labs, museums,
and rapid prototyping, where the bottleneck is modelling skill rather than
printer access.

\paragraph{Problem statement.}
The promise does not survive contact with a 3D printer, for two reasons.

\begin{enumerate}[leftmargin=*]
  \item \textbf{Hardware access.} TRELLIS.2 and its dependencies are written
        for CUDA. A fabrication lab equipped with Apple Silicon machines
        cannot run it at all.
  \item \textbf{Generated meshes are rendering assets, not manufacturing
        assets.} A generative model is trained to look correct from the
        outside. Nothing in its objective requires the mesh to be watertight,
        manifold, or free of disconnected debris --- and in practice it is
        none of these. Our test asset is a single photographed object, but
        the generator emits it as \textbf{4{,}518 disconnected components},
        not watertight, with a hollow interior. A slicer given this file
        either refuses it or silently produces a nonsense toolpath.
\end{enumerate}

Repairing the mesh is possible, but the repair is where the real difficulty
lies. The only reliable way to \emph{guarantee} watertightness is to throw the
original topology away and rebuild the surface from a voxel grid --- an
operation that reshapes the object. The user is therefore asked to accept a
trade: manufacturability in exchange for fidelity. Existing tools ask the user
to accept that trade \emph{blind}. They report success or failure, never the
magnitude of the distortion they introduced.

\paragraph{Key assumptions and scope.}
We assume a single, roughly object-centric input photograph and consumer
fused-deposition printing. We assume the generated mesh, however malformed, is
geometrically \emph{correct enough} --- our validation layer measures deviation
introduced by \emph{post-processing}, not the accuracy of the generation itself
(Section~\ref{sec:limitations}). We do not attempt to auto-correct thin walls
or overhangs, because both require shape-altering edits that we consider unsafe
to apply without user consent.

\paragraph{Proposed system.}
We contribute a three-stage pipeline:

\begin{itemize}[leftmargin=*]
  \item \textbf{Generation.} A port of the TRELLIS.2 4B flow-matching model
        from CUDA to PyTorch~MPS with Metal backends, producing a textured mesh
        from one image in $\approx$5~minutes on an M4~Pro.
  \item \textbf{Print preparation.} Noise-component filtering, best-effort
        repair with a guaranteed voxel-remesh fallback, exterior flood-fill to
        make the object solid, appearance re-baking, and decimation.
  \item \textbf{A validation layer.} A measurement stage that reports, for
        every run, how far the printable mesh has moved from the mesh the
        generator actually produced, together with printability diagnostics.
\end{itemize}

The validation layer is the contribution we consider most important, and it is
the focus of this report. It is what converts an opaque, irreversible repair
into an auditable one.

% ===========================================================================
\section{Related Work}
% ===========================================================================

\paragraph{Generative image-to-3D.}
TRELLIS \cite{xiang2025trellis} introduced Structured LATents (SLat), coupling
a sparse voxel grid with dense multi-view features from a vision foundation
model, decodable to radiance fields, Gaussians, or meshes. TRELLIS.2
\cite{xiang2025trellis2} replaces SLat with \emph{O-Voxel}, a field-free sparse
voxel structure that encodes geometry and PBR appearance jointly and --- key
for our use case --- supports arbitrary topology including open, non-manifold,
and fully-enclosed surfaces. We adopt TRELLIS.2 as our generator because it is
the strongest available single-image model with native PBR output. Its
conditioning stack supplies our other dependencies: DINOv3
\cite{simeoni2025dinov3} for image features and RMBG-2.0, built on BiRefNet
\cite{zheng2024birefnet}, for background removal.

Crucially, none of these works target manufacturability. TRELLIS.2's ability to
represent non-manifold and open surfaces is a \emph{feature} for rendering and
precisely the \emph{problem} for printing: the representation is free to emit
geometry a printer cannot interpret. This is the gap our work addresses.

\paragraph{Direct mesh generation.}
An alternative to repairing dense output is generating clean topology in the
first place. TreeMeshGPT \cite{lionar2025treemeshgpt} autoregressively produces
compact, artist-like meshes with orders of magnitude fewer faces. We evaluated
it as a replacement for the TRELLIS.2 decoder and did not adopt it: it
generates untextured geometry from a point cloud rather than PBR assets from an
image, and its output is optimised for editability, not watertightness --- so
it would not remove our need for a repair-and-validate stage. We report it here
because the comparison motivates our design: no current generator outputs
print-ready topology, so \emph{some} validated repair stage is unavoidable.

\paragraph{Mesh repair and measurement.}
Our repair path rests on classical geometry: marching cubes
\cite{lorensen1987marchingcubes} to rebuild a surface from the filled voxel
grid, and quadric error metric decimation \cite{garland1997qem}, via
\texttt{fast\_simplification} \cite{fastsimplification}, to bring the inflated
face count back down. Our fidelity metric follows the standard symmetric
point-to-surface formulation of Aspert et al.~\cite{aspert2002mesh}, whose
central argument --- that mesh error must be measured surface-to-surface rather
than vertex-to-vertex when topology differs --- is exactly our situation.
Implementation uses trimesh \cite{trimesh} and xatlas \cite{xatlas}. The gap we
fill is one of integration rather than algorithms: these error metrics are
standard in mesh-processing research but are absent from generative
image-to-3D tooling, which reports repair as a boolean rather than a
measurement.

% ===========================================================================
\section{Methodology}
% ===========================================================================

\subsection{System overview}

The pipeline runs in three stages, each a standalone command:

\begin{lstlisting}[caption={The pipeline end to end.}]
# Stage 1: photo -> textured mesh (TRELLIS.2 on MPS)
python generate.py photo.png

# Stages 2 and 3: mesh -> printable mesh, validated
python make_printable.py output_3d.glb
\end{lstlisting}

\subsection{Stage 1: Generation on Apple Silicon}

Porting a 4B-parameter CUDA pipeline to MPS is largely a matter of surviving
the places where the Metal stack behaves unlike CUDA. Three were significant:

\begin{itemize}[leftmargin=*]
  \item \textbf{Backend selection with fallbacks.} Sparse convolution runs on
        a Metal \texttt{flex\_gemm} backend where available, falling back to a
        pure-PyTorch path when the \texttt{metallib} fails to load (it ships
        MSL~4.0, which only loads on macOS~26+). Attention uses SDPA.
        \texttt{PYTORCH\_ENABLE\_MPS\_FALLBACK} must be set before torch is
        imported anywhere, or unsupported ops crash rather than falling back
        to CPU.
  \item \textbf{The macOS GPU watchdog.} Long-running Metal kernels in the
        SLat decoder are killed by the OS watchdog. The failure is
        pernicious: it does \emph{not} raise a Python exception. Execution
        continues with empty tensors and fails later with an unrelated error.
        We detect the known downstream signatures and report the real cause
        with actionable workarounds instead of a misleading stack trace.
  \item \textbf{Texture baking.} A Metal-accelerated baker is used when
        available, with a pure-Python KDTree baker (UV-unwrapped via xatlas
        \cite{xatlas}) as fallback.
\end{itemize}

\subsection{Stage 2: Print preparation}
\label{sec:prep}

Given the generated \texttt{.glb}, the preparation stage applies:

\begin{enumerate}[leftmargin=*]
  \item \textbf{Component filtering.} The mesh is split into connected
        components and ranked by bounding-box diagonal. Components below
        \texttt{--min-object-ratio} (default 5\%) of the largest are discarded
        as floating-dot noise. This runs in every mode, since generated assets
        carry thousands of specks that a slicer would treat as real objects.
  \item \textbf{Best-effort repair.} Vertex merging, degenerate-face removal,
        hole filling, and normal fixing. If this achieves watertightness and
        solid infill is not requested, the pipeline stops here --- this is the
        \emph{fidelity-preserving} path.
  \item \textbf{Guaranteed repair (voxel remesh).} Otherwise the mesh is
        voxelised and the interior flood-filled, then a surface is rebuilt with
        marching cubes \cite{lorensen1987marchingcubes}. Two details matter:
        \begin{itemize}
          \item Filling uses the \texttt{orthographic} method, not the default
                hole-filling method. The latter fills only fully-enclosed
                voids, so any pinhole in the voxel shell lets the fill leak and
                the object stays a hollow shell --- which appears in the slicer
                as floating islands. The orthographic method fills every voxel
                between the first and last surface voxel along each axis,
                producing one solid mass even when the shell leaks.
          \item Voxelisation uses \texttt{max\_iter=None} so the subdivision
                depth is computed from the mesh's longest edge. The library
                default caps at 10 passes and simply fails on fine pitches.
        \end{itemize}
  \item \textbf{Appearance re-baking.} The voxel path destroys the original
        topology and UV atlas, so colour is resampled onto the new vertices by
        closest-point and barycentric interpolation from the original surface.
  \item \textbf{Decimation.} Marching cubes inflates the face count, so the
        mesh is decimated back toward a target using QEM \cite{garland1997qem}.
        The collapse map returned by \texttt{fast\_simplification} is replayed
        to keep per-vertex colours consistent with the decimated vertices.
\end{enumerate}

Objects retain their original orientation. An earlier version auto-oriented
each object onto a ``best'' resting pose by convex-hull search; we removed it
because it silently moved the user's model and the choice of pose is better
left to the slicer.

\subsection{Stage 3: The validation layer}
\label{sec:validation}

Every step in Section~\ref{sec:prep} is lossy, irreversible, and invisible. The
validation layer exists to make the loss visible. It has two halves: a
\emph{fidelity report} that measures deviation from the generator's output, and
a \emph{printability report} that flags manufacturing risks.

\subsubsection{Why not compare vertices?}

The obvious way to compare a mesh before and after repair --- match vertices
and measure displacement --- is unusable here, because repair does not preserve
correspondence. It does not even preserve \emph{count}: in our measurements the
voxel path takes the 1.1M-vertex, 2.23M-face input, rebuilds an entirely new
surface, and decimates it to a different mesh at 45.5\% of the original face
count (Table~\ref{tab:results}). Every vertex is new; there is no vertex $i$ in
the output corresponding to vertex $i$ in the input.

We therefore measure \textbf{surface to surface}, following
\cite{aspert2002mesh}. We sample $n = 20{,}000$ points uniformly by area from
each mesh and compute point-to-surface distances in both directions. This
requires only that both meshes occupy the same coordinate frame --- which holds
because our pipeline never moves the object --- and is completely agnostic to
vertex count, ordering, and topology.

\subsubsection{Fidelity metrics}

Let $A$ be the pre-repair mesh, $B$ the processed mesh, and $S_A, S_B$ the
sampled point sets. With $d(p, M)$ the distance from point $p$ to the closest
point on the surface of mesh $M$, we report:

\begin{equation}
d_{\mathrm{chamfer}} =
  \frac{1}{2}\!\left(
  \frac{1}{|S_A|}\sum_{p \in S_A} d(p, B) +
  \frac{1}{|S_B|}\sum_{q \in S_B} d(q, A)
  \right)
\label{eq:chamfer}
\end{equation}

\begin{equation}
d_{\mathrm{hausdorff}} = \max\!\left(
  \max_{p \in S_A} d(p, B),\;
  \max_{q \in S_B} d(q, A)
  \right)
\label{eq:hausdorff}
\end{equation}

Both are reported in absolute units \emph{and} as a percentage of the model's
bounding-box diagonal, which makes them scale-free and comparable across
objects of different sizes.

The two metrics are deliberately complementary, and reporting only one would be
misleading:

\begin{itemize}[leftmargin=*]
  \item \textbf{Chamfer} (Eq.~\ref{eq:chamfer}) is a mean, so it describes
        typical deviation --- but it \emph{averages away} localised damage. A
        repair that is perfect everywhere except that it amputates one thin
        antenna still reports an excellent Chamfer score.
  \item \textbf{Hausdorff} (Eq.~\ref{eq:hausdorff}) is a maximum, so it catches
        exactly that amputation. It is correspondingly sensitive to single
        outliers.
\end{itemize}

Read together, a low Chamfer with a high Hausdorff is the signature of
\emph{localised} damage --- the case that matters most, because it is the case
a user would notice on the printed object and that a mean-only metric hides.

Two further quantities are reported:

\begin{itemize}[leftmargin=*]
  \item \textbf{Volume change (\%)}, reported \emph{only} when both meshes are
        watertight. Volume is undefined for a non-watertight mesh, and since
        the input to repair is typically not watertight, this field is
        usually \texttt{N/A}. Guarding it rather than printing a meaningless
        number is deliberate: a validation layer that emits confident nonsense
        is worse than one that admits ignorance.
  \item \textbf{Face count ratio (\%)}, which exposes the topological cost of
        the repair.
\end{itemize}

\subsubsection{Printability diagnostics}

The second half of the layer flags manufacturing risks, and \emph{reports
rather than fixes} them:

\begin{itemize}[leftmargin=*]
  \item \textbf{Overhang area.} The fraction of total surface area whose face
        normal points downward beyond the 45$^\circ$ rule, i.e.\ faces with
        $n_z < -\cos(45^\circ)$, weighted by face area. Predicts support
        material requirements.
  \item \textbf{Thin walls.} 500 rays are cast inward from randomly sampled
        face centres along the inward normal; the first-hit distance
        approximates local wall thickness. Distances below 1.0\,mm
        ($\approx 2.5\times$ a typical 0.4\,mm nozzle) are counted as warnings.
\end{itemize}

Neither is auto-corrected. Thickening a wall or adding a chamfer changes the
object's visible shape, and we judged silent shape edits a worse failure than
an informed warning --- consistent with the reasoning that removed
auto-orientation.

\subsection{Evaluation Procedure}
\label{sec:eval}

\paragraph{Tools.} PyTorch (MPS backend), TRELLIS.2-4B \cite{xiang2025trellis2},
trimesh 4.12.2 \cite{trimesh}, scikit-image (marching cubes),
\texttt{fast\_simplification} \cite{fastsimplification}, xatlas \cite{xatlas}.
Hardware: Apple M4~Pro, 24\,GB unified memory.

\paragraph{Test input.} We evaluate on a real asset produced by Stage~1 from a
single photograph --- not a synthetic or hand-authored mesh, since the
pathologies we target are precisely those a generator introduces. Its
properties are given in Table~\ref{tab:input}.

\paragraph{Baseline.} Defining a baseline requires care, because the naive
comparison (``mesh before vs.\ after'') is what our metric already computes. We
therefore compare the two \emph{policies} the tool can apply, which represent
the two genuine options a user has:

\begin{itemize}[leftmargin=*]
  \item \textbf{Baseline --- best-effort repair} (\texttt{--no-solid-infill}):
        conventional cleanup, preserves topology, but offers \emph{no
        watertightness guarantee}.
  \item \textbf{Full system --- guaranteed solid} (default): voxel remesh with
        exterior flood fill; always watertight and solid.
\end{itemize}

\paragraph{Metrics.} Watertightness (the binary gate that decides whether the
object is printable at all), Chamfer and Hausdorff deviation, volume change,
face ratio, overhang area, thin-wall warnings, and wall-clock time.

\paragraph{Correctness check on the metric itself.} A validation layer must
itself be validated. We verify that when post-processing is a genuine no-op
(an already-watertight mesh through the best-effort path), the fidelity report
returns exactly zero on every metric --- confirming that non-zero readings
reflect real geometric change and not sampling noise in the measurement.

% ===========================================================================
\section{Results}
% ===========================================================================

\begin{table}[h]
\centering
\caption{The generated test asset, as emitted by TRELLIS.2 from one
photograph. This is the input to print preparation.}
\label{tab:input}
\begin{tabular}{lr}
\toprule
\textbf{Property} & \textbf{Value} \\
\midrule
Vertices                     & 1{,}113{,}494 \\
Faces                        & 2{,}229{,}064 \\
Watertight                   & \textcolor{kaorange}{\textbf{No}} \\
Connected components         & \textcolor{kaorange}{\textbf{4{,}518}} \\
Components kept as real geometry & 1 \\
\bottomrule
\end{tabular}
\end{table}

\noindent
Table~\ref{tab:input} is itself a result, and the headline one: a photograph of
\emph{one} object becomes a mesh of \textbf{4{,}518 disconnected pieces}, of
which exactly one is the object and 4{,}517 are noise. This is the concrete
form of the gap identified in Section~1 --- the generator's output is
visually convincing and structurally unusable.

\begin{table}[h]
\centering
\caption{Full print-preparation run on the asset in Table~\ref{tab:input},
default (guaranteed-solid) policy. Every number is measured by the validation
layer on a single run; timing is wall-clock on an M4~Pro.}
\label{tab:results}
\begin{tabular}{lrr}
\toprule
\textbf{Quantity} & \textbf{Input} & \textbf{Processed} \\
\midrule
Watertight (printable at all)        & \textcolor{kaorange}{No}  & \textcolor{kalime}{\textbf{Yes}} \\
Faces                                & 2{,}229{,}064 & $\approx$1{,}000{,}000 \\
Face count vs.\ input surface        & 100\%         & 45.5\% \\
\midrule
\multicolumn{3}{l}{\emph{Fidelity of processed vs.\ input surface (Eqs.~\ref{eq:chamfer},~\ref{eq:hausdorff})}}\\
Chamfer distance (mean deviation)    & ---           & 0.00100 \;(\textbf{0.09\%}) \\
Hausdorff distance (worst deviation) & ---           & 0.05033 \;(\textbf{4.68\%}) \\
Volume change                        & ---           & N/A (input not watertight) \\
\midrule
\multicolumn{3}{l}{\emph{Printability diagnostics}}\\
Overhang area ($>45^\circ$)          & ---           & 31.2\% of surface \\
Thin-wall warnings ($<1$\,mm, 500 rays) & ---        & 499 \\
\midrule
Repair time                          & ---           & 78\,s \\
Total wall-clock                     & ---           & 104\,s \\
\bottomrule
\end{tabular}
\end{table}

\noindent
The default and \texttt{--no-solid-infill} policies produced an
\emph{identical} result on this asset (both 87\,s repair, same output mesh),
because best-effort cleanup alone could not make the mesh watertight --- so both
fell through to the voxel path. This is itself a finding
(Section~\ref{sec:discussion}): on real generative output the ``preserve
topology'' baseline is not actually available, and the guaranteed path is not
optional.

\begin{table}[h]
\centering
\caption{Validation-layer self-check: an already-watertight mesh through the
best-effort path. Exact zeros confirm the metric reports real geometric change
rather than sampling noise.}
\label{tab:selfcheck}
\begin{tabular}{lr}
\toprule
\textbf{Metric} & \textbf{Value} \\
\midrule
Chamfer distance   & 0.00000 (0.00\%) \\
Hausdorff distance & 0.00000 (0.00\%) \\
Volume change      & +0.00\% \\
Face count ratio   & 100.0\% \\
\bottomrule
\end{tabular}
\end{table}

% ===========================================================================
\section{Discussion}
\label{sec:discussion}
% ===========================================================================

\paragraph{The two metrics disagree, and that is the result.}
The single most informative number in Table~\ref{tab:results} is not either
distance on its own but their \emph{ratio}. The Chamfer distance is
0.09\% of the model size --- by any reasonable standard, an excellent repair:
averaged over the whole surface, the printable mesh sits a tenth of a percent
away from what the generator produced. Read alone, it would justify shipping the
result without a second thought. The Hausdorff distance is 4.68\%, roughly
\textbf{50 times larger}. This is exactly the localised-damage signature
predicted in Section~\ref{sec:validation}: the repair is faithful almost
everywhere and badly wrong in a small region --- most likely a thin protrusion
that the voxel grid could not resolve and marching cubes rounded off. A
mean-only report would have hidden precisely the defect a user would notice on
the printed object. This is the concrete payoff of reporting both metrics, and
it is why we consider the validation layer, not the port, the contribution.

\paragraph{The guards fire correctly.}
Two fields report \texttt{N/A} or a self-check zero rather than a number, and
both are correct to do so. Volume change is \texttt{N/A} because the input mesh
is not watertight, so its volume is genuinely undefined --- printing a
plausible-looking percentage there would be fabrication. And on the identity
self-check (Table~\ref{tab:selfcheck}) every metric reads exactly zero,
confirming that the 0.09\%/4.68\% readings above are real geometric change and
not an artifact of area sampling. A measurement tool earns trust by being
verifiably silent when there is nothing to report; ours is.

\paragraph{An unexpected finding: the baseline collapsed.}
We designed the evaluation around two policies, expecting the best-effort path
to preserve topology at some fidelity cost and the solid path to guarantee
printability at a larger one. On real generative output the distinction
vanished: best-effort repair could not close the mesh, so both policies took the
voxel path and returned the same result. The lesson is that conventional mesh
cleanup, which assumes near-manifold input, is simply the wrong tool for the
output of an image-to-3D model --- the geometry is too far from watertight for
hole-filling to recover. This validates the decision to make the aggressive
voxel path the default rather than an opt-in.

\paragraph{Where the diagnostics meet, and miss, expectations.}
The overhang figure (31.2\%) is plausible for an organic, curved object and is
the kind of number a user can act on by adding supports. The thin-wall count
(499 of 500 rays) is the opposite --- it is almost certainly an artifact, not a
measurement, because the mesh is in normalised coordinates where the fixed
1\,mm threshold flags essentially the entire surface. We report it honestly as a
limitation (Section~\ref{sec:limitations}) rather than dressing it up: a
diagnostic that fires on everything conveys no information, and the fix is a
physical scale, not a different threshold.

% ===========================================================================
\section{Limitations and Future Work}
\label{sec:limitations}
% ===========================================================================

\begin{itemize}[leftmargin=*]
  \item \textbf{We validate the repair, not the reconstruction.} This is the
        most important limitation of the work. Our metrics measure deviation
        between the processed mesh and \emph{the mesh the generator produced}
        --- not between the processed mesh and the real photographed object. A
        hallucinated feature that TRELLIS.2 invents is faithfully preserved and
        reported as 0\% deviation. The pipeline therefore certifies
        \emph{``we did not damage what the model gave us,''} which is strictly
        weaker than \emph{``this is an accurate replica.''}
        \emph{Next step:} photograph the printed object from known viewpoints
        and measure against the input image, or evaluate against a
        ground-truth scanned dataset to characterise generation error
        separately from repair error.
  \item \textbf{The thin-wall metric assumes millimetre units.} The 1.0\,mm
        threshold is compared directly against mesh coordinates, so it is only
        meaningful if the model is scaled to millimetres. Generated assets are
        typically in an arbitrary normalised frame, which makes the reported
        warning count unreliable in the default pipeline.
        \emph{Next step:} require an explicit physical scale (or target print
        size) and convert before thresholding; report the metric as
        \texttt{N/A} when no scale is known, following the same guard already
        used for volume change.
  \item \textbf{Thin-wall sampling is sparse and stochastic.} 500 rays over a
        multi-hundred-thousand-face mesh is a thin sample, and it is unseeded,
        so counts vary between runs on identical input.
        \emph{Next step:} seed the sampler for reproducibility and scale the
        ray budget with surface area; validate against a slicer's own
        thickness analysis.
  \item \textbf{Single test object.} All quantitative results come from one
        generated asset. The numbers characterise a case, not a distribution.
        \emph{Next step:} run the pipeline over a benchmark set spanning
        thin-featured, multi-part, and organic shapes, and report metric
        distributions rather than point values.
  \item \textbf{No end-to-end print verification.} We validate geometry, but
        we have not closed the loop by printing the validated meshes and
        confirming that the diagnostics predict real print outcomes.
        \emph{Next step:} print a sample spanning the range of reported
        overhang and thin-wall scores and correlate the predictions with
        observed failures --- the only way to establish that the thresholds
        are calibrated rather than plausible.
  \item \textbf{Hausdorff is unbounded by construction.} A single stray voxel
        speck surviving the noise filter dominates the worst-case metric.
        \emph{Next step:} additionally report a robust quantile (e.g.\ 95th
        percentile) alongside the true maximum to separate systematic
        deviation from isolated outliers.
\end{itemize}

% ===========================================================================
\section{Conclusion}
% ===========================================================================

We built a single-image-to-3D-print pipeline that runs on Apple Silicon and,
unlike existing image-to-3D tooling, measures the geometric cost of the repair
it performs. A photograph becomes a textured mesh in about five minutes via a
CUDA-to-MPS port of TRELLIS.2; that mesh --- 2.2M triangles across 4{,}518
disconnected pieces, not watertight --- is then made printable by a repair stage
whose every step is lossy and irreversible. The core contribution is the
validation layer that renders those losses visible: using topology-agnostic
symmetric Chamfer and Hausdorff surface distances, it certified our test asset
as average-faithful (0.09\% mean deviation) while flagging a localised worst-case
error 50 times larger (4.68\%) that a mean-only metric would have concealed, and
it reads exactly 0.00\% when post-processing changes nothing --- proving the
measurement reports real change rather than noise. The takeaway is a shift in
what an image-to-3D tool is accountable for: not merely ``did it produce a
printable file,'' but ``how far did making it printable move it from what the
model generated, and where'' --- an answer the user can finally see before they
commit the object to a printer.

\printbibliography

\end{document}
